% Ejemplo ilustrativo (cifras de resultados hipotéticas). Sustitúyelo por el texto de tu investigación.

\guia{La introducción representa entre 8 y 12\% del documento y sigue este orden. Las etiquetas grises señalan cada elemento; desaparecen al ocultar las guías.}

\etiqueta{1. Planteamiento del problema}
El peso mexicano es una de las monedas de economías emergentes con mayor volumen de negociación en el mercado cambiario global \parencite{bis2022}; sin embargo, el volumen anual de contratos de futuros negociados en el Mercado Mexicano de Derivados (MexDer) cayó alrededor de 86\% entre 2009 y 2022, aun en episodios de alta volatilidad cambiaria como las elecciones estadounidenses de 2016 o la pandemia de COVID-19 en 2020. Esta paradoja es relevante porque un mercado local de derivados profundo y líquido permite a empresas, bancos e inversionistas cubrir su exposición cambiaria a menor costo y con menor riesgo de contraparte.

\etiqueta{2. Contexto y antecedentes}
La teoría de la cobertura sostiene que los agentes con exposición en el mercado de contado toman posiciones en futuros para reducir la varianza de su riqueza \parencite{johnson1960}, y que la efectividad de esa cobertura puede medirse como la proporción de la varianza que se elimina \parencite{ederington1979}. Bajo este enfoque, una mayor volatilidad debería traducirse en mayor demanda de cobertura y, por lo tanto, en mayor volumen negociado. No obstante, la evidencia para México se ha concentrado en la eficiencia de cobertura de los contratos del MexDer y ha prestado poca atención a la competencia de mercados externos, como el CME Group, y del mercado extrabursátil, lo que deja un vacío para explicar la contracción del mercado local después de 2015.

\etiqueta{3. Objetivo e hipótesis}
Esta investigación tiene como objetivo analizar la relación entre la volatilidad del tipo de cambio peso-dólar y el volumen de negociación de contratos de futuros del dólar en el MexDer durante el período 2009–2022. La hipótesis de trabajo sostiene que la volatilidad cambiaria incrementa el volumen de futuros del dólar en el MexDer, pero que este efecto se ha debilitado conforme la cobertura migra hacia el CME y hacia el mercado extrabursátil.

\etiqueta{4. Metodología general}
La investigación adopta un enfoque cuantitativo de alcance correlacional-explicativo. Se construye una base de datos de series de tiempo mensuales (2009–2022) con información del Banco de México, el MexDer y el CME Group. La volatilidad condicional del tipo de cambio se estima con un modelo GARCH(1,1) y su efecto sobre el volumen negociado se evalúa con regresiones lineales con errores estándar robustos (Newey-West), estimadas para el período completo y para los subperíodos 2009–2015 y 2016–2022.

\etiqueta{5. Variables}
La variable dependiente es el logaritmo del volumen mensual de contratos de futuros del dólar negociados en el MexDer. Las variables independientes son la volatilidad condicional del tipo de cambio FIX, el diferencial entre la TIIE y la tasa de fondos federales de Estados Unidos y el logaritmo del volumen de futuros del peso en el CME; como variables de control se incluyen el índice VIX y una variable dicotómica para la pandemia de COVID-19.

\etiqueta{6. Estructura del documento}
El documento se organiza en tres capítulos y un apartado de conclusiones. El Capítulo 1 presenta el marco teórico sobre la cobertura con futuros; el Capítulo 2 desarrolla el marco contextual, con la evolución del mercado de derivados en México; el Capítulo 3 expone la metodología y el análisis de resultados, que incluye la descripción de los hallazgos, su contraste con la literatura y sus implicaciones de política pública y regulación; finalmente, las conclusiones retoman los hallazgos, las limitaciones y las líneas futuras de investigación.
